\chapter{Fundamentals}\label{ch:fundamentals}

This chapter establishes the fundamentals required for the rest of the thesis: the cache hierarchy
and the notion of cache awareness, trie-based index structures with the Adaptive Radix Tree as a
recurring reference, the unified \texttt{std::map} comparison interface as a search-algorithm shell
with standard interfaces, the unified \texttt{std::vector} comparison interface as a container
shell, the workload frameworks used in the literature (YCSB and others such as SOSD, TPC, SPEC~CPU,
CloudSuite, and mimalloc-bench), the quality criteria and patterns of scientific measurement, and
the C++23 metaprogramming facilities on which the cache engine builds.

\section{Cache Hierarchy and Cache Awareness}\label{sec:cache-basics}
\subsection{Memory and Cache Levels}\label{ssec:cache-levels}
Modern processors interpose several cache levels (L1, L2, and --- not on every platform --- L3)
between the registers and main memory~\cite{hennessy2019architecture}. The transfer unit between
levels is the \emph{cache line}; its size, however, is not a universal constant but a property fixed
per microarchitecture (Section~\ref{ssec:cache-line-sizes}). Nor does transfer occur strictly
``block-wise in one piece'': techniques such as \emph{critical-word-first} deliver the requested
word first; DDR5 splits the memory channel into two independent 32-bit
sub-channels~\cite{micron_ddr5}; and non-temporal or write-combining stores bypass the normal
write-allocate path rather than fetching the affected cache line into the cache
hierarchy~\cite{intel_sdm}. An access that finds its datum in a cache level is cheap; a miss is
forwarded to the next, slower level, and a miss in the last-level cache reaches main memory at a
cost of several hundred cycles~\cite{hennessy2019architecture,drepper2007memory}. On multi-socket
and chiplet systems, non-uniform memory access (NUMA) additionally makes the latency of a reference
depend on which node owns the page~\cite{drepper2007memory}.

\subsection{Cache-Line Sizes: Static and Dynamic}\label{ssec:cache-line-sizes}
The cache-line size is hard-wired per microarchitecture but differs between architectures. On x86-64
(AMD as well as Intel) and on common ARM cores (e.g.\ Cortex-A76 or Neoverse~V2) it is 64~bytes;
individual designs use wider lines at the SoC level --- Apple Silicon, for instance, 128~bytes on L2
and the system-level cache~\cite{apple_optimization}. Across architectures, the cache-line sizes in
use range from narrower 32~bytes on older or embedded cores through the now-common 64-byte line up
to 128~bytes~\cite{hennessy2019architecture,intel_optimization,amd_sog}. Programs can \emph{read}
the actual size at run time --- via \texttt{CPUID} on x86, the \texttt{CTR\_EL0} register on ARM,
and --- with limitations --- the \texttt{riscv\_hwprobe} interface on RISC-V (e.g.\ CBO block
sizes)~\cite{arm_arm,intel_sdm} --- but cannot change it. The depth of the hierarchy likewise
varies: not every platform has an L3 cache (simple RISC-V cores such as the SiFive~U74 stop at a
shared L2), while others bring a large or even asymmetrically distributed L3 (e.g.\ the stacked 3D
V-Cache of some AMD processors). For a cache-line-aware layout, the line size and cache depth
determined per target platform are therefore an \emph{input parameter of compilation} rather than a
fixed assumption; Section~\ref{ssec:aware-oblivious} takes this up again for the delimitation of the
thesis.

\subsection{Address Translation: TLB and dTLB}\label{ssec:tlb}
Besides the data-cache hierarchy there exists a second one, for address translation: the
\emph{Translation Lookaside Buffer} (TLB) caches virtual-to-physical mappings, and a TLB miss
triggers an expensive page walk~\cite{hennessy2019architecture}. For locality reasons, processors
separate the \emph{instruction TLB} (iTLB) from the \emph{data TLB} (dTLB) --- since instruction
accesses are markedly more local than data accesses ---; a multi-level arrangement (L1 dTLB plus a
shared L2 STLB) is also common~\cite{intel_sdm,drepper2007memory}. For search structures, the dTLB
is the second, independent locality boundary alongside the cache line: large or widely scattered
nodes raise dTLB pressure independently of cache utilization, which is why huge pages (2~MiB/1~GiB)
enlarge the TLB reach.

\subsection{Cache Coherence, Write and Update Behaviour per Implementation}\label{ssec:coherence}
Caches handle writes according to two basic strategies: \emph{write-through} writes immediately
through to the next level, whereas \emph{write-back} initially records the change only in the cache
line (\emph{dirty}) and writes it back only on eviction --- saving bus traffic and lowering write
latency~\cite{hennessy2019architecture}. Orthogonally, the \emph{allocation policy} determines how a
write miss is handled: \emph{write-allocate} first loads the target line into the cache (typical
with write-back), \emph{no-write-allocate} writes straight through to the next level (typical with
write-through). ARMv8 makes this choice explicit through page attributes --- write-through,
write-back, or non-cacheable, as well as read- or read-write-allocate~\cite{arm_arm}, while x86 uses
write-back with write-allocate as the normal case~\cite{intel_sdm}.

Which line is \emph{evicted} and thereby replaced on a miss is decided by a
\emph{heuristic replacement strategy}. Besides the exact and approximate classics (LRU, pseudo-LRU),
scan- and thrash-resistant methods are common: \emph{Re-Reference Interval Prediction} (RRIP,
dynamic DRRIP) makes do with two bits per line~\cite{jaleel2010rrip} and builds on the adaptive
insertion policy~\cite{qureshi2007adaptive}. How many levels may hold a line at once is governed,
finally, by the \emph{inclusion policy}: an \emph{inclusive} cache keeps every line of the higher
levels redundantly, an \emph{exclusive} (victim) cache excludes them, and a \emph{non-inclusive} one
allows both~\cite{hennessy2019architecture}.

In multi-core systems a \emph{coherence protocol} ensures the consistency of shared lines; the
widespread, invalidation-based \emph{MESI} protocol keeps each line in one of the states
\emph{Modified}, \emph{Exclusive}, \emph{Shared}, or \emph{Invalid}, extended by vendors into
variants such as MOESI (AMD) or MESIF (Intel)~\cite{hennessy2019architecture}. Which write,
allocation, replacement, inclusion, and coherence strategy applies concretely is predominantly a
property of the respective microarchitecture, cache level, and the memory attributes set by the
operating system, and not freely selectable at run time; their concrete manifestations are presented
by the state of the art (Section~\ref{sec:sota-cache}), and their effect --- in particular
\emph{cache-line ping-pong} under contending writes to the same line --- motivates the telemetry and
concurrency axes of the engine (Chapters~\ref{ch:concept} and~\ref{ch:impl}).

\subsection{Cache-Aware and Cache-Oblivious}\label{ssec:aware-oblivious}
A data structure is \emph{cache-aware} if its layout is tuned to concrete cache parameters (line
size, level capacities), and \emph{cache-oblivious} if it achieves good locality across all levels
without knowing them~\cite{frigo1999cache}. For tree-shaped structures the cache-oblivious
arrangement is theoretically grounded: Bender, Demaine, and Farach-Colton minimise the expected
number of block transfers when laying out a tree of fixed topology to within an additive constant of
optimal~\cite{bender2002treelayout} and maintain dynamic ordered sets with locality-preserving
$O(k/B)$ traversal~\cite{bender2002scanning}. Both philosophies recur throughout the state of the
art (Chapter~\ref{ch:sota}) and motivate the layout, prefetch, and hardware axes of the proposed
engine.

The present thesis pursues a third, intermediate path. \emph{First}, the actual cache properties
(line size, hierarchy depth, topology) are \emph{determined} per target platform and enter a
\emph{system-optimized binary} as fixed quantities --- the line size itself is a compilation input,
since on the platforms considered here it is not changeable at run time as a program parameter
(Section~\ref{ssec:cache-line-sizes}). \emph{Second}, the \emph{configurable} cache-near settings
--- prefetch distance, batch and inline thresholds, pool budgets, and, where architecture and
operating system permit, the hardware prefetcher --- are carried as axis parameters that are
permuted at run time against an already-loaded binary (Section~\ref{ssec:dynamics-levels} and
Chapter~\ref{ch:concept}). \emph{Third}, this latitude is delimited depending on hardware and
operating system: from fully run-time-configurable (mutable Linux with MSR access) through boot- or
compile-time static (immutable operating systems) to fixed (read-only cache registers on ARM and
RISC-V). Thus cache-oblivious behaviour is complemented by \emph{measured} locality knowledge that
enables a heuristic, measurement- oriented selection of the settings per load pattern; the
\emph{why} of this construction is developed in Chapter~\ref{ch:concept}.

\subsection{Heuristics and Their Automation}\label{ssec:heuristics}
The ``measurement-oriented selection of the settings'' hinted at in the previous section rests on
the notion of a \emph{heuristic}. In the context of cache optimization a heuristic is a
\emph{decision rule without a guarantee of optimality}: it maps observable features --- such as the
cache properties determined per target platform (Section~\ref{ssec:cache-line-sizes}) or the
measured load profile (Section~\ref{sec:workloads-basics}) --- to a concrete configuration without
proving that configuration best. It takes the place of an exact solution wherever the latter is
unknown or too expensive to determine; the cache's own heuristic replacement strategies (LRU, RRIP;
Section~\ref{ssec:coherence}) are the hardware-side archetype of this way of
thinking~\cite{jaleel2010rrip}.

Heuristics can be distinguished by \emph{when} their decision is made --- this corresponds to the
levels of dynamics developed in Section~\ref{ssec:dynamics-levels}. A \emph{static heuristic} is
fixed at compile time: the chosen configuration is compiled firmly into a system-optimized binary
and no longer changes at run time. A \emph{dynamic heuristic}, by contrast, selects the
configuration at run time, depending on the actually observed behaviour, and may adjust it during
operation. Both variants presuppose the same \emph{heuristic function} but shift the moment of its
evaluation.

Formally, this heuristic function is a \emph{mapping} from a measurement result or load profile to a
configuration: $h: \mathcal{M} \rightarrow \mathcal{K}$, where $\mathcal{M}$ denotes the space of
collected features (axis measurements, load-profile metrics) and $\mathcal{K}$ the space of
admissible axis assignments. The present thesis first \emph{measures} this relationship
systematically (Chapter~\ref{ch:methodology}); the \emph{automated} determination of $h$ itself ---
that is, the mechanism that derives the respective best configuration from the measurement results,
towards an ML classification of load-profile and axis features onto configurations --- is an
objective and outlook of this thesis (Section~\ref{sec:outlook}) and not yet realized.

Finally, the terms used partly synonymously in the literature must be delimited. A heuristic is
called \emph{adaptive} if it adjusts its decision to changing conditions; \emph{dynamic} refers more
narrowly to the moment --- evaluation at run time rather than at compile time; and
\emph{measurement-driven} emphasizes the data source --- the decision rests on collected
measurements rather than on fixed model assumptions. An adaptive heuristic is thus always dynamic, a
dynamic one not necessarily adaptive, and both may, but need not, be measurement-driven.

\section{Classes of Search Structures}\label{sec:search-classes}
Search structures can be classified by the mechanism with which they map a key to a
record~\cite{knuth1998taocp3}. \emph{Comparison-based} structures locate the key via order
comparisons and presuppose a total order; these include the balanced search trees from AVL and
red-black trees to the block-oriented B-/B\textsuperscript{+} trees. \emph{Digital} structures
branch by the bytes of the key rather than by comparison (tries, radix trees, ART). \emph{Hashing}
maps keys directly via a hash function and has no inherent order. \emph{Spatial} structures,
finally, partition a multi-dimensional space (k-d tree, quadtree, R-tree). This thesis concentrates
on the \emph{tree-like} structures of the comparison- and the digital-based class --- the common
denominator of modern in-memory indexes ---, while hashing appears only as a comparison baseline and
spatial structures lie outside the scope. The principle developed in Chapter~\ref{ch:concept} is,
however, not bound to trees: through a higher-level abstraction, further structure classes can be
attached without rebuilding the measurement and permutation apparatus (Section~\ref{sec:outlook})
--- for however different these classes appear on the outside, they share a common inner
\emph{anatomy} of orthogonal partial decisions that Chapter~\ref{ch:concept} exposes as an axis set.

Within these mechanism classes, research has produced concrete structure \emph{classes}; they are
named and substantiated here at the class level, with one section per mechanism class, while their
internal dissection into axes --- including concrete internal manifestations such as ART's node
types --- occurs only in the state of the art (Chapter~\ref{ch:sota}) as \emph{instances}.

\subsection{Comparison-Based Search Trees}\label{ssec:comparison-trees}
Comparison-based trees locate keys via order comparisons and, when balanced, attain logarithmic
search depth; they differ mainly in how they maintain this balance and how large their nodes are.

\subsubsection{AVL and Red-Black Trees}
From the unbalanced binary search tree arise the self-balancing
\emph{AVL trees}~\cite{adelsonvelsky1962avl}, whose sibling subtrees differ in height by at most
one, and the \emph{red-black trees}~\cite{bayer1972symmetric,guibas1978redblack}, which secure the
same logarithmic height through a node colouring with looser but cheaper-to-maintain invariants.

\subsubsection{B- and B\texorpdfstring{\textsuperscript{+}}{+} Trees}
For block- and external-memory-oriented indexes, the \emph{B-trees}~\cite{bayer1972btree} and the
\emph{B\textsuperscript{+} trees}~\cite{comer1979btree} became established; they keep all payload in
the leaves and reduce inner nodes to signposts, so that one node typically matches a memory or
cache-block size~\cite{graefe2001btree,hankins2003nodesize}.

\subsubsection{T-Tree}
Directly pertinent to main memory is the \emph{T-tree}~\cite{lehman1986ttree}, which combines the
balancing of the AVL tree with node-local value ranges and thereby lowers the number of pointer
indirections per search step.

\subsection{Digital and Prefix-Based Structures}\label{ssec:digital}
Digital structures branch by the bytes of the key rather than by comparison, so that the search-path
length depends on the key length and not on the number of keys.

\subsubsection{Trie}
The \emph{trie}~\cite{delabriandais1959trie,fredkin1960trie} is the basic form: every path from the
root spells out a key, and common prefixes are shared between keys.

\subsubsection{Radix and PATRICIA Trees}
\emph{Radix} and \emph{PATRICIA} trees~\cite{morrison1968patricia} collapse sparsely populated
single-child chains through path compression and thereby lower height and memory footprint compared
with the plain trie.

\subsubsection{Adaptive Radix Tree (ART)}
The \emph{Adaptive Radix Tree} (ART)~\cite{leis2013art} adapts the node representation to the actual
branching degree and serves this thesis throughout as a reference design; \emph{fanout} and
\emph{node representation} thereby determine both space requirement and cache behaviour.

\subsection{Hashing}\label{ssec:hashing}
Hashing maps keys directly via a hash function and has no inherent order, which is why range queries
are expensive~\cite{knuth1998taocp3}. Besides open and chained addressing,
\emph{cuckoo hashing}~\cite{pagh2004cuckoo} achieves constant worst-case lookups, and modern,
SIMD-accelerated tables such as SwissTable~\cite{abseil_swisstable} attain high practical density.
Hashing appears in this thesis only as a non-tree-like comparison baseline.

\subsection{Spatial Structures}\label{ssec:spatial}
Spatial structures partition a multi-dimensional key space --- e.g.\ the
\emph{k-d tree}~\cite{bentley1975kdtree}, the \emph{quadtree}~\cite{finkel1974quadtree}, and the
\emph{R-tree}~\cite{guttman1984rtree}. With their multi-dimensional keys they lie outside the scope
of this thesis and are mentioned only for completeness.

\subsection{Flat Methods and Theoretical Exotics}\label{ssec:flat-exotic}
As a limiting case stands the \emph{flat} method: a sorted array with binary or interpolation
search~\cite{knuth1998taocp3} is, strictly speaking, the single leaf of a tree. On string keys the
\emph{suffix array} generalises this flat, sorted layout into a space-efficient full-text index
whose lightweight, memory-budgeted construction is a model example of algorithm
engineering~\cite{manzini2002suffixarray}. Theoretical ``exotics'' shift the complexity bounds ---
the \emph{van Emde Boas tree}~\cite{vanemdeboas1975} and the
\emph{y-fast trie}~\cite{willard1983yfast} achieve $O(\log\log u)$ over a key universe of size $u$,
and the \emph{fusion tree}~\cite{fredman1993fusion} undercuts the information-theoretic comparison
bound. They are theoretically significant but, for the cache-line-aware in-memory index of this
thesis, only a marginal consideration.

\subsection{Subdivision of Containers}\label{ssec:container-taxonomy}
Orthogonal to the search-algorithm view stands the established subdivision of the \emph{containers}
of the C++ standard library~\cite{iso_cpp}: \emph{sequence containers} (e.g.\ \texttt{std::vector})
order elements by insertion position, \emph{ordered associative} containers (\texttt{std::map},
\texttt{std::set}) by key order, and \emph{unordered} variants (\texttt{std::unordered\_map}) map
via a hash function. These containers provide standardized interface contracts against which
structurally different methods can be measured uniformly (Section~\ref{sec:compare-interfaces}): the
ordered map is the common denominator for comparison- and digital-based search structures, a
\texttt{std::vector} provides the sequence view, and a hash table
(e.g.\ \texttt{absl::flat\_hash\_map}/SwissTable) the unordered counter-check.

\subsection{Conceptual Placement in the Present Framework}\label{ssec:own-framework}
The present thesis places these external methods into a framework of its own that
Chapter~\ref{ch:concept} develops in detail and that is only introduced here. At the topmost level
stands the \emph{genus} as an externally visible interface --- SearchAlgorithm, Container, and
Graph; beneath it lie \emph{living-being subclasses} with a fixed axis set (the SearchAlgorithm
subclass with its nineteen axes, as well as Set, Sequence, Adapter, and View under the Container
interface). An \emph{axis} is an orthogonal partial decision --- an ``organ'' --- of an overall
algorithm; a concrete method is thus a point in the Cartesian product of the axes, and the classical
families of Sections~\ref{ssec:comparison-trees}--\ref{ssec:flat-exotic} turn out to be mere reading
groupings of orthogonal axis assignments~\cite{idreos2018periodic,idreos2018datacalculator}. As its
own probe this thesis contributes PRT-ART (a Probst Redirect Tree / ART hybrid), whose axis
contributions are likewise dissected only in Chapters~\ref{ch:concept} and~\ref{ch:impl}. Thus the
external classification (Sections~\ref{ssec:comparison-trees}--\ref{ssec:container-taxonomy}) and
the axis-oriented dissection view used in this thesis remain cleanly separated.

\section{Unified Comparison Interfaces}\label{sec:compare-interfaces}
\subsection{The \texttt{std::map} Interface}\label{sec:stdmap-interface}
To compare structurally different search algorithms fairly, they are all mapped onto a single,
well-understood interface: the C++ contract of the ordered map \texttt{std::map<Key,\,Value>}
(\texttt{find}/\texttt{insert}/\texttt{erase}/range iteration). The axes introduced in
Chapter~\ref{ch:concept} describe the \emph{inner} behaviour behind this interface, while the
interface itself stays fixed --- the comparison of two conformant \texttt{std::map} implementations
is thus the central measurement act. A variadic specialization keeps the interface ergonomic: one
template parameter provides a \texttt{std::vector}-like API, two parameters the \texttt{std::map}
API, and $N>2$ a \texttt{std::map<K,\,std::tuple<\ldots>>}. The core contract of the ordered map
comprises the point query (\texttt{find}, \texttt{at}, \texttt{operator[]}), insertion and update
(\texttt{insert}, \texttt{insert\_or\_assign}, \texttt{emplace}), deletion (\texttt{erase}), and
order-based navigation (\texttt{lower\_bound}, \texttt{upper\_bound}, range
iteration)~\cite{iso_cpp}; the complete listing of \emph{all} standard functions with expected
semantics and corresponding conformance tests is given in Appendix~\ref{app:interfaces}.

\subsection{The \texttt{std::vector} Interface}\label{ssec:stdvector-interface}
Parallel to the map shell, the sequence shell \texttt{std::vector} serves as the container
comparison case~\cite{iso_cpp}. Its one-parameter contract offers indexed access
(\texttt{operator[]}, \texttt{at}), resizing (\texttt{push\_back}, \texttt{emplace\_back},
\texttt{insert}, \texttt{erase}, \texttt{resize}), and contiguous storage (\texttt{data}); shared
with the map are iteration and \texttt{size}/\texttt{empty}/\texttt{clear}, while the addressing
differs (position rather than key order). This makes it possible to check whether an axis composite
also carries container semantics and not only search semantics; the complete function table of both
shells is in Appendix~\ref{app:interfaces}.

\subsection{Levels of Dynamics}\label{ssec:dynamics-levels}
Behind the fixed external interface the system distinguishes several \emph{levels of dynamics} that
determine \emph{when} an axis decision is made:
\begin{itemize}
  \item \textbf{Compile time vs.\ run time} as the topmost distinction;
  \item the \textbf{living-being type}, which occupies only a certain subset of the available axes;
  \item the \textbf{axis algorithms}, which are compiled into a living-being binary at compile time
  (static recombination);
  \item the \textbf{dynamic sub-axes}, whose run-time parameters (prefetch distance, batch size,
  inline threshold, pool budget) are varied against an already-loaded binary without recompilation;
  \item the \textbf{static sub-axes} as a compile-time selection among static partial algorithms of
  differing complexity.
\end{itemize}
This separation is anchored in the permutation tree of the cache engine (Chapter~\ref{ch:concept}):
static axis levels each produce their own binary per path, dynamic sub-axes are virtual run-time
loops over the same binary. Only thereby does it become explainable, per interface function, how it
is internally wired --- an aspect that a dedicated cache-engine developer document deepens. Relative
to the state of the art --- the static design-space enumeration of the
\emph{Data Calculator}~\cite{idreos2018datacalculator} and classical
\emph{policy-based design}~\cite{alexandrescu2001modern} --- this taxonomy carries compile-time- and
run-time-variable axis levels jointly and explicitly; a web search found no established naming for
this layering (cf.\ Section~\ref{ssec:new-patterns}).

\section{Loads and Workloads}\label{sec:workloads-basics}
\subsection{Term: Load, Workload, and Goal of Variation}\label{ssec:load-terms}
A \emph{load} (operation) is a single call of the search interface --- a point query, an insertion,
a deletion, or a range scan. A \emph{workload} is a structured profile of such operations: an
operation mix (e.g.\ 50\,\% read / 50\,\% update) over a key access distribution (uniform,
Zipf-distributed, latest values) on a dataset of defined size and key characteristic. Realistic
workloads are thereby \emph{playbooks} of many consecutive interface calls --- not isolated,
unrealistic single accesses.

The \emph{goal of variation} is to break a methodological bias: each publication typically chooses
the load profile under which its own algorithm performs best (a read-only trie avoids updates, a
compressed filter seeks negative queries, a cache-aware layout favours Zipf hotspots). Only when
\emph{all} load profiles run against \emph{all} living-being binaries (and additionally the
container shells) can largely bias-reduced comparability be gained. The concrete ``all-against-all''
matrix belongs in the evaluation methodology (Chapter~\ref{ch:methodology}); here only the
conceptual framework is laid.

\subsection{Workload Frameworks}\label{sec:ycsb}
The workloads used in the literature stem from \emph{several} benchmark frameworks that this thesis
considers jointly; the most prominent is the Yahoo!\ Cloud Serving Benchmark
(YCSB)~\cite{cooper2010ycsb}. It defines a family of key-value workloads that serve here as a common
load profile: YCSB-A (50/50 read/update, Zipf-distributed), YCSB-B (95/5 read/update), YCSB-C (read
only), YCSB-D (latest values), YCSB-E (short range scans), and YCSB-F (read-modify-write). Each
algorithm profile declares an expected workload (Chapter~\ref{ch:sota}) that the measurement driver
uses for profile-aware routing (Chapter~\ref{ch:methodology}).

Alongside this, the analysed publications rely on further frameworks: the index benchmark
\emph{SOSD}~\cite{kipf2019sosd}, the transaction benchmarks of the \emph{TPC} family (TPC-C,
TPC-DS)~\cite{tpc_benchmarks}, integration into real systems (e.g.\ SuRF~\cite{zhang2018surf} as a
Bloom-filter replacement in RocksDB), real string corpora (URLs, DNA, dictionary terms), and --- for
the hardware-near works --- \emph{SPEC~CPU}~\cite{spec_cpu2017} and
\emph{CloudSuite}~\cite{ferdman2012cloudsuite}; the allocator corpus is measured with
allocator-specific suites (\texttt{mimalloc-bench}~\cite{mimalloc_bench}, Larson, threadtest,
shbench, and LADDIS). Because these frameworks produce different, mutually incomparable loads, the
thesis abstracts them into a common set of load profiles and runs these as a
\emph{dynamic workload axis} over all living-being binaries (Chapter~\ref{ch:methodology}); the
connection of the non-YCSB frameworks occurs via a dataset loader.

\section{Scientific Measurement}\label{sec:measurement-basics}

\subsection{Quality Criteria of Scientific Measurement}\label{ssec:measurement-criteria}
Scientific measurement is subject to established quality criteria~\cite{hoefler2015benchmarking}.
\emph{Reproducibility} requires fixed random seeds, a deterministic execution order, and fully
documented measurement configurations, so that runs are methodically repeatable and statistically
comparable. \emph{Validity} breaks down into internal (causally cleanly isolated effects), construct
(the metric actually measures the intended concept), and external validity (transferability to other
platforms and loads). \emph{Fairness} demands identical inputs --- the same keys, workloads, and
warmup --- for all comparison candidates; missing capabilities of a baseline are reported as
\emph{N/A} rather than covertly emulated through auxiliary constructs. Finally, \emph{dispersion} is
part of the statement: instead of arithmetic means, percentiles over HDR histograms are reported and
the statistical significance of differences is tested~\cite{hoefler2015benchmarking}.

\subsection{Measurement Patterns and Concepts Used}\label{ssec:measurement-patterns}
On these criteria several measurement patterns build, whose \emph{why} only Chapter~\ref{ch:concept}
develops. A \emph{two-phase operation loop} separates a discarded warmup run from the actual
measurement and thus decouples latency from the accumulated state. Collection occurs \emph{hybrid}
via two paths: isolated axis algorithms are measured directly against one another, while a composite
living being is collected centrally via an ABI-stable observer interface (per-axis observer). A
\emph{conformance gate} checks before each measurement that an implementation satisfies the
\texttt{std::map} contract, so that only comparable things are compared; each configuration is
collected in several separate runs with reported dispersion~\cite{hoefler2015benchmarking}.

These concepts form the \emph{definitional classes} of measurement; the state of the art
(Chapter~\ref{ch:sota}) presents the \emph{instances} that technically realize them ---
e.g.\ hardware performance counters (PMC), the OTF2/VAMPIR profiling tools, or the TU-Dresden
hardware measurement methodology. Thus fundamentals and state of the art mesh directly.

\section{C++23 Metaprogramming and Design Patterns}\label{sec:cpp23}
\subsection{Classical Design Patterns}\label{ssec:classic-patterns}
The system uses named design patterns (GoF)~\cite{gamma1994patterns} throughout; only those actually
used in the code are introduced here. \emph{Strategy} encapsulates each of the nineteen axes as an
interchangeable algorithm; \emph{Abstract Factory} creates the static and dynamic nodes of the
permutation tree, \emph{Composite} forms that tree, and an \emph{Iterator} traverses it lazily.
\emph{Adapter} binds foreign code behind the stable engine interface (Habich directive).
\emph{Observer} (pull model) and \emph{Decorator} (non-intrusive measurement shells around the axes)
collect the measurements, and a \emph{Memento} (copy-on-write) saves and discards the state between
warmup and measurement. Static polymorphism arises through the
\emph{Curiously Recurring Template Pattern} (CRTP) together with \emph{Template Method} concept
checks. Not used, by contrast, are Singleton, Command, Flyweight, and Bridge.

\subsection{Metaprogramming Facilities}\label{ssec:metaprog-facilities}
The engine relies on compile-time composition to avoid run-time dispatch in the hot path.
\emph{Concepts} constrain the building-block interfaces; the
\emph{Curiously Recurring Template Pattern} (CRTP) provides static polymorphism;
\texttt{if~constexpr} and \texttt{std::conditional\_t} choose per axis between a probe
implementation and a cache-engine default; and C++23 \emph{modules} structure the translation units
between the permutation binaries and the builder, but do not by themselves guarantee portable ABI
stability~\cite{wg21_modules_abi}; that is established only by an explicitly C-stable module
boundary (\texttt{comdare\_create\_anatomy}). Together these facilities realize the
\emph{zero-cost abstraction} on which the permutation matrix rests. Added to this are
\emph{policy-based design} (the composition as a tuple of nineteen axis policies),
\emph{tag dispatch} for the probe-versus-default resolution, and \emph{type-list computation}
(Boost.MP11), which describes the axis product space without fully materializing
it~\cite{alexandrescu2001modern}.

\subsection{New and Unnamed Metaprogramming Patterns}\label{ssec:new-patterns}
Beyond the established patterns, two constructions occur for which a web search found no established
naming and which are therefore named as a contribution of this thesis. First, the
\emph{lazy-sparse Cartesian materialization}: the product space of all axis variants is only counted
arithmetically at compile time --- over $10^{14}$ possible axis combinations, of which the subset
actually materialized as standalone binaries is correspondingly smaller --- and materialized
\emph{individually} per path, mixed-radix and on demand, instead of instantiating the full type tree
(which overwhelmed the compiler). Second, \emph{type-driven source emission}: from a compile-time
axis composition the fully qualified type names are extracted, rendered to C++ source, compiled per
permutation as a standalone module, and loaded across the ABI boundary --- so the downloaded probe
(PRT-ART) is compiled via metaprogramming into permutation/matrix layers against the core
constituents of the cache engine. Both patterns combine known idioms (concepts, CRTP, type lists)
into an architecture-specific novelty; their full realization is developed in
Chapters~\ref{ch:concept} and~\ref{ch:impl}.
